\scanpage{170}
\mspage{44}

a) $\unsure{\mathcal M|M}$ stable par petites limites inductives (et pas
seulement filtrantes)

b) $\mathcal M$ \deleted{``accessible'' (\unsure{objectivement}), i.e.} muni d'une
filtration cardinale $(\mathrm{Filt}^{\unsure{\pi|\varpi}})_{\pi\geq\pi_0}$

c) Les objets de $\mathcal M$ sont accessibles (NB. il suffit
de \unsure{le|la} supposer pour les objets de $\mathrm{Filt}^{\pi_0}$).

b) implique que $\mathcal M$ contient une petite
famille génératrice par épim. stricts (p. ex.
une sous-catégorie \inserted{$C$} petite de $\underline{\mathrm{Filt}}^{\pi_0}$ \inserted{$=\mathcal M^{\pi_0}$} telle
que l'inclusion $C\unsure{-|\to}\underline{\mathrm{Filt}}^{\pi_0}$ soit une
équivalence de catégories). Par SGA 4 I 8.12.7,
il en résulte que $\mathcal M$ est stable par petites
limites projectives. Il s'ensuit le

\uline{Lemme 5.1.} (Sous les conditions a), b) ) Pour
tout objet $X$ de $\mathcal M$, l'ensemble $\mathrm{Ss}(X)$
des sous-objets de $X$, muni de la relation
d'ordre naturelle ``d'inclusion'', \deleted{\unsure{\}|;}} admet des
$\mathrm{Inf}$ et des $\mathrm{Sup}$ quelconques

Des $\mathrm{Inf}$, car ceux-ci s'expriment comme
des limites projectives dans $\mathcal M$. Or l'existence
des $\mathrm{Inf}$ implique celle des $\mathrm{Sup}$,
puisque $\underset{J}{\mathrm{Sup}}=\underset{J'}{\mathrm{Inf}}$, où $J'$ est l'ens.
des majorants de $J$ dans $\mathrm{Ss}(X)$.

\uline{Cor. 5.2} : Pour \unsure{toute|\deleted{tout}} flèche $f:Z\longrightarrow X$
dans $\mathcal M$, il existe, parmi les sous-objets \inserted{$X'$} de $X$
par lesquels $f$ se factorise, un plus petit

% archive sheet number: 169
